\begin{table*}[t]
\centering
\small
\caption{\textbf{Example of self-revision.}
An A2 candidate that fails shortcut resistance (L8) and its revision. The gold answer is B in both versions.}
\label{tab:example_selfrev}
\begin{tabular}{@{}p{0.07\textwidth}p{0.43\textwidth}p{0.43\textwidth}@{}}
\toprule
 & \textbf{Korean} & \textbf{English translation} \\
\midrule
Before
& Q-system으로 터널 지보량을 검토할 때 ESR을 선정해야 한다. 다음 지하 암반시설들 중 굴착지보비 값을 가장 크게 잡는 것이 타당한 것은?
\newline A) 지하 원자력 발전소
\newline B) 채굴 작업 구간 옆에 마련되어 장비 대기와 버력 적치에 쓰이다가 해당 구간 채굴이 끝나면 되메워질 공동
\newline C) 지하 식품저장 시설
\newline D) 수력발전소의 수로터널
& The ESR must be selected when reviewing tunnel support with the Q-system. For which of the following underground rock facilities is it appropriate to set the largest excavation support ratio?
\newline A) An underground nuclear power plant
\newline B) A cavern prepared beside a mining section, used for equipment standby and muck storage, to be backfilled once mining of the section ends
\newline C) An underground food storage facility
\newline D) The headrace tunnel of a hydropower plant \\
\midrule
Audit (L8)
& \multicolumn{2}{p{0.87\textwidth}@{}}{\textsc{fail}: The gold is isolated by form. B is the only long, detailed situational option, while A, C, and D are short generic facility labels.} \\
\midrule
After
& Q-system으로 터널 지보량을 검토하기 위해 굴착지보비(ESR)를 선정하려고 한다. 다음 지하 암반시설들 중 ESR 값을 가장 크게 잡는 것이 타당한 것은?
\newline A) 지하 원자력 발전소에서 운전 설비를 계속 수용하도록 계획된 주설비 공동
\newline B) 채굴 작업 구간 옆에서 장비와 버력을 두었다가 작업 종료 뒤 되메우도록 계획된 보조 공동
\newline C) 지하 식품저장 시설에서 식품 재고를 계속 수용하도록 계획된 저장 공동
\newline D) 수력발전소에서 발전용 용수를 계속 통수하도록 계획된 수로터널
& The excavation support ratio (ESR) is to be selected in order to review tunnel support with the Q-system. For which of the following underground rock facilities is it appropriate to set the largest ESR?
\newline A) A main cavern in an underground nuclear power plant, planned to keep housing operating equipment
\newline B) An auxiliary cavern beside a mining section, planned to hold equipment and muck and to be backfilled after the work ends
\newline C) A storage cavern in an underground food storage facility, planned to keep housing food stock
\newline D) A headrace tunnel in a hydropower plant, planned to keep conveying water for power generation \\
\bottomrule
\end{tabular}
\end{table*}